\section{Falsifiability conditions}
\label{sec:falsify}

The claims of this paper are falsified if at least one of the following
conditions holds.
Each condition specifies a concrete way in which the formal structure
asserted here would fail.

\begin{enumerate}
  \item \textbf{Absence of threshold-induced discontinuities.}
  If all phase predicates defined in \artifact{ETS.K_EA.v1.1} can be
  reformulated as continuous, precedence-free conditions such that
  phase classification varies continuously with respect to admissible
  parametrizations of states in $\Omega(K_{EA})$, then the proposed
  analogy to threshold-induced phase transitions is invalid.

  \item \textbf{Failure of boundary correspondence under faithful mapping.}
  If empirical or simulated enterprise trajectories, when mapped
  consistently and without reinterpretation to ETS variables, do not
  exhibit boundary-like changes in phase membership at predicate sign
  changes or loss of observability, then the diagnostic claim of
  threshold-induced phase reassignment is refuted.

  \item \textbf{Violation of absorbing or no-go constraints.}
  If \texttt{STOP} or any other ETS-defined no-go invariant admits
  outgoing phase-level transitions, or behaves as non-absorbing under
  the ETS allowance relation, then the formal phase logic asserted in
  this paper is invalid.
\end{enumerate}

All falsifiability conditions are stated at the level of formal structure.
They do not presuppose any specific empirical dataset, but they admit
operational testing wherever a faithful, non-substitutive mapping to
ETS variables is available.
